\section{Realising the Structure: Global Maps, Local Transport, and Composition}

\subsection{Prediction, transport, and composition}

Three tasks, graded separately throughout: \textbf{prediction} --- from a fixed anchor, forecast features at other angles; \textbf{transport} --- apply a specified increment from a starting state never used in fitting; \textbf{composition} --- apply several operations whose product must equal one true total transformation, graded against the ground-truth composite target as well as for internal consistency. Each has its own holdouts (Appendix C). The three tasks are not three benchmarks but the three faces of Definition 1 weighed separately: prediction tests approximation along the orbit, transport tests transfer to states never fitted, and composition tests the homomorphism law $\rho$($\tau$ $_2$ $\tau$ $_1$) = $\rho$($\tau$ $_2$)$\circ$ $\rho$($\tau$ $_1$) itself.

\subsection{Global realisations}

Five natural global families were fitted at matched budgets, and exactly one of them combines fidelity with composition: the band-limited family predicts features of genuinely re-rendered inputs to \textbf{0.17--0.40\% relative MSE}, up to $16.6\times$ tighter than the copy baseline. Content conditioning does not rescue global forms (48 matched-capacity contrasts, including an oracle partition by the true displacement direction); \S{}4.6 states what that negative result does and does not imply about the cross-content reading of the residual error at measured sites. \textbf{Table 7} gives the numbers.

One of the rejected families fails \textbf{by theorem rather than by observation}.

\textbf{Theorem 3 (a closed orbit has no exact coordinate-monotone realisation).} A family $\{G_t\}$ whose members are coordinate-wise monotone in a fixed basis and whose dependence on the family parameter is monotone --- $G_t(z)_i=f_{i,t}(z_i)$ with $(t,u)\mapsto f_{i,t}(u)$ non-decreasing in both arguments --- realises an orbit only if every coordinate of that orbit is monotone in the parameter. Hence a closed non-degenerate orbit admits no exact realisation by such a family, and each coordinate that decreases by $A$ over intervals of length $\ell$ contributes $L^2$ error at least $A^2\ell/4$. (Proved in Appendix B.5.1; the two lemmas are B.7 and B.8.)

The theorem says what the class cannot do; the measurement says what it costs. On the reference orbit the best coordinate-monotone fit carries $L^2$ error $998$ against $924$ for the best fixed linear map (Table 7), and both fail by a wide margin, while the carrier family that keeps the conjugate pair is exact. Segmenting the basis --- letting it follow the field locally, \S{}5.3 --- is the escape the theorem permits and the measurements confirm.


\begin{table*}[t]
\centering
\small
\caption*{\textbf{Table 7.} Global realisations at matched budget, each against the copy floor (\S{}3.7).}
\begin{tabularx}{\textwidth}{@{}>{\raggedright\arraybackslash}X >{\raggedright\arraybackslash}X >{\raggedright\arraybackslash}X >{\raggedright\arraybackslash}X@{}}
\toprule
\textbf{family} & \textbf{fidelity vs copy} & \textbf{composition} & \textbf{verdict} \\
\midrule
symmetric & $0.09$--$0.64\times$ & fails & rejected \\
antisymmetric & $3.07\times$ & fails & rejected \\
band-limited multi-frequency & $0.1$--$0.6\times$ & fails & rejected \\
low-order discrete ($\mathbb Z_2/\mathbb Z_3$) & $\le 0.41$ & fails & rejected \\
monotone per-dimension & no exact realisation of a closed orbit (Thm 3); best fit measured, $L^2=998$ on the reference orbit & fails by theorem and by measurement & rejected \\
\textbf{band-limited, $K=3$--$4$} & \textbf{$0.17$--$0.40\%$ relative MSE; $6.7$--$16.6\times$ over copy} & passes & \textbf{the working global family} \\
content-conditioned (48 matched contrasts) & mean gain $-0.167$, best $+0.017$ & --- & structurally ineffective, not merely weak \\
\bottomrule
\end{tabularx}
\end{table*}

Segmenting --- letting the basis follow the field locally --- is exactly the escape this theorem permits, which is the local transporter of \S{}5.3.

\subsection{Local state-conditioned realisations}

The successful realisation is \textbf{local and phase-conditioned}: the hue circle is divided into $40^\circ$ segments, each carrying a generator estimated from paired endpoint features, with the current phase read from the representation and the segment generator applied --- a nonlinear, state-dependent composite. It fits an order of magnitude tighter than the best global family and closes loops to $0.4\%$ on unseen shapes and starting states. \textbf{Table 8} collects the comparison. Two further gaps --- the median advantage of a local step over a global single direction at a fixed site, and how that gap grows with depth --- are not reported here, because no producing artefact exists for them.


\begin{table*}[t]
\centering
\small
\caption*{\textbf{Table 8.} Local versus global realisation at matched budget.}
\begin{tabularx}{\textwidth}{@{}>{\raggedright\arraybackslash}X >{\raggedright\arraybackslash}X >{\raggedright\arraybackslash}X@{}}
\toprule
\textbf{quantity} & \textbf{value} & \textbf{} \\
\midrule
local vs copy, per site & $1.55$--$15.94\times$ &  \\
global vs copy, same sites & $0.08$--$0.84\times$ &  \\
generator: third-order Fourier fit in phase & $R^2 = 0.87$ & method box, Appendix E \\
orbit bending rate across layers & $1.91\to2.35$ & Appendix D \\
\bottomrule
\end{tabularx}
\end{table*}

\textit{One mechanism hypothesis fits this evidence: the obstacle to global linearisation is shape--colour resonance rather than curvature or non-sharing. Leave-one-shape-out linearisation succeeds $4/4$ at the mid-level site and fails $2/4$ at the deep site, exactly where phase organisation destabilises --- the signature a resonance account predicts. It is stated as a hypothesis, not a result. Figure 6 contrasts the two realisations site by site and prices the composition of a fixed total transformation.}

\begin{figure*}[t]
\centering
\includegraphics[width=\textwidth]{fig06_realisation.pdf}
\caption*{\textbf{Figure 6.} Which realisations work, and what composing costs. (a) Gain over the copy baseline of the best local phase-conditioned transporter (circles) and of the best global family (squares) at five sites, logarithmic axis, with a dotted line at one (no better than copying); the local construction is $1.55$--$15.9\times$ the copy baseline while the global family is $0.08$--$0.84\times$. (b) Transfer at a fixed total hue rotation as a function of how the rotation is split into a chain, with the direct fit as a dashed reference; bars are means over three seeds and whiskers one standard deviation. (c) The same for the heat semigroup, whose chains degrade monotonically with length.}
\end{figure*}

\subsection{Composable transport and its price}

Composition is where the algebra is tested. Holding the \textbf{total} transformation fixed and varying only its decomposition --- two- and four-step paths, an order swap, a cyclic round trip that must return to the identity, and, for heat, chains of two, four and eight steps against a single-step substitute control and a direct fit as the upper reference --- chains come within a few points of a direct fit for both algebras while the single-step substitute does not, and heat degrades monotonically in chain length. Path \textit{consistency} and path \textit{fidelity} come apart: the global linear family agrees most tightly between paths and transports worst, so a family of operators can be perfectly self-consistent and still fail to implement the transformation. \textbf{Table 9} names every path and carries every number, including the saturation caveat (the heat single step scores high because that total sits near saturation) and the cyclic round-trip measure.


\begin{table*}[t]
\centering
\small
\caption*{\textbf{Table 9.} Composition at a fixed total transformation (hue $90^\circ$; heat $t^\star=1.0$), plain-CNN stage 1, three seeds, zero split overlap between fit and evaluation.}
\begin{tabularx}{\textwidth}{@{}>{\raggedright\arraybackslash}X >{\raggedright\arraybackslash}X >{\raggedright\arraybackslash}X >{\raggedright\arraybackslash}X@{}}
\toprule
\textbf{total} & \textbf{path} & \textbf{ridge $T_F$} & \textbf{conv-res $T_F$} \\
\midrule
hue $90^\circ$ & direct fit at $90^\circ$ & 0.285 $\pm$ 0.008 & 0.511 $\pm$ 0.019 \\
hue $90^\circ$ & $45+45$ & 0.237 $\pm$ 0.013 & 0.472 $\pm$ 0.023 \\
hue $90^\circ$ & $30+60$ & 0.249 $\pm$ 0.012 & 0.488 $\pm$ 0.027 \\
hue $90^\circ$ & $60+30$ (order swap) & 0.245 $\pm$ 0.014 & 0.484 $\pm$ 0.028 \\
hue $90^\circ$ & $22.5\times4$ & 0.253 $\pm$ 0.024 & 0.451 $\pm$ 0.032 \\
hue $90^\circ$ & single $22.5^\circ$ ($1/4$ of total) & 0.131 $\pm$ 0.019 & 0.207 $\pm$ 0.010 \\
hue $90^\circ$ & cyclic $+120/-120$ (identity) & 0.067 $\pm$ 0.029 & 0.254 $\pm$ 0.021 \\
heat $t^\star=1.0$ & direct fit at $t^\star$ & 0.610 $\pm$ 0.107 & 0.760 $\pm$ 0.030 \\
heat & two steps & 0.514 $\pm$ 0.132 & 0.679 $\pm$ 0.019 \\
heat & four steps & 0.437 $\pm$ 0.112 & 0.625 $\pm$ 0.063 \\
heat & eight steps & 0.372 $\pm$ 0.114 & 0.481 $\pm$ 0.025 \\
heat & single step ($t^\star/8$) & 0.572 $\pm$ 0.025 & 0.597 $\pm$ 0.025 \\
\bottomrule
\end{tabularx}
\end{table*}

One law organises the whole table, and it is a statement about the \textit{truth} action rather than about any operator.

\textbf{Theorem 4 (accumulation law).} Let the truth action $\tau$ be $L$-Lipschitz in the consumer's seminorm, and let a realised family have single-step error $E_1=\varepsilon$ uniformly on the states visited. Then the $n$-step chain error satisfies
\[\adjustbox{max width=\columnwidth}{$\displaystyle \begin{aligned}
E_n&\le\varepsilon+L\,E_{n-1},\qquad\text{equivalently}\\
E_n&\le\varepsilon\frac{L^n-1}{L-1}\ (L\neq1),\qquad E_n\le n\varepsilon\ (L=1).
\end{aligned}
$}\]
\textit{Proof.} Insert and subtract the truth's intermediate state: $\rho(\tau)^n z-g_{\tau^n}z=\big[\rho(\tau)(\rho(\tau)^{n-1}z)-g_\tau(\rho(\tau)^{n-1}z)\big]+\big[g_\tau(\rho(\tau)^{n-1}z)-g_\tau(g_{\tau^{n-1}}z)\big]$; the first bracket is at most $\varepsilon$ by the uniform single-step bound, the second at most $L\,E_{n-1}$ by Lipschitzness. Iterating gives the closed forms. Appendix B carries the full statement and the uniformity hypotheses. $\square$

The bound separates three regimes by one property of the truth: $L<1$ (contractive) gives a \textbf{bounded} geometric accumulation $\varepsilon/(1-L)$; $L=1$ (isometric) gives \textbf{linear} accumulation $n\varepsilon$; $L>1$ (expansive) gives \textbf{exponential} growth $L^n$. The two tested families differ in how their chains behave, and the accumulation law is the tool that separates the possible causes: a contraction caps the chain error, an isometry accumulates it linearly, an expansion grows it exponentially. Which regime a site occupies must be \textit{measured} in the consumer's seminorm, not inferred from the generator. Two readings follow, and they are worth separating. Heat's monotone degradation in Table 9 is therefore \textbf{not} amplification by the truth action --- a contraction forbids that --- but a growing single-step error as the chain wanders into states where the fitted family is worse; and hue's flat chains, under the loosest of the three bounds, mean that its single-step error is already small rather than that the bound protects it. The regime that would actually blow up, $L>1$, is not instantiated in this paper. In all three cases the composition law itself is automatic; what the accumulation law prices is the realisation error. Estimating $\varepsilon$ and $L$ independently of the measured chain error, and instantiating the expansive class ($L>1$), where the bound predicts exponential growth, are not done here.

\textbf{Metric distortion at real sites.} The one number Theorem 4 needs is how the truth action distorts distances in the consumer's seminorm, and it is measurable. Over pairs of held-out crops, ResNet-50, four depths and both algebras, the median ratio of feature distances after and before the transformation is below one for heat and near one for hue in the \textbf{site} metric --- the qualitative prediction of the generators --- but the same ratio in the \textbf{consumer's} response sits slightly above one for both algebras at every depth (Table D.22). An input-space contraction is not a contraction of what a downstream consumer reads. That is why Theorem 4 is stated with $L$ in the consumer's seminorm and why a site's regime is reported as a distribution of typical stretch rather than read off the generator.

\subsection{Functional consequences of the realised action}

Capability claims are graded on ground-truth targets and downstream behaviour, because the proxies can be gamed. The random orthogonal operator of \S{}1.3 attains the best projection score at $15$ of $42$ sites while its transfer is negative at every one of them and its detector agreement is $F_1=0.000$. Shuffled-pair controls and PCA-aligned subspace controls complete the picture, and the proxies are not worthless: within a family, projection correlates with transfer ($\rho=+0.93$), and the feature-residual ratio correlates with the correct sign --- each instrument is informative inside its scope and insufficient alone. The evaluation rule this paper uses follows directly: feature-target fidelity, composition behaviour, and downstream effect, each against its own baseline, jointly.

Realisation capability varies across layers, and the variation is structured. The next section shows it follows the readout--transformation pairing, and what learning does to it.

